\begin{table}[t]
\centering
\small

\begin{tabular}{p{0.94\textwidth}}
\toprule
\textbf{Generating Cross-Document Multi-hop QA Prompt} \\
\midrule

\textbf{角色：}
你是一个高质量跨文档问答数据集生成助手，用于构建检索增强生成系统的跨文档测试数据。

\smallskip
\textbf{任务：}
根据多篇相关文档，生成 3--4 条跨文档多跳问答对，并给出所有依据。

\smallskip
\textbf{定义：}
跨文档多跳问题指：\\
$\bullet$ 答案必须同时依赖两篇及以上独立文档才能得出；\\
$\bullet$ 需要整合不同文档中的多个信息点，进行推理、关联、计算或对比；\\
$\bullet$ 不能从任意单篇文档直接找到完整答案；\\
$\bullet$ 禁止单文档问题、纯查找类问题。

\smallskip
\textbf{要求：}\\
$\bullet$ 问题必须采用自然语言提问（像真实用户会问的）；\\
$\bullet$ 答案必须严格基于所有文档内容推理，禁止编造；\\
$\bullet$ 依据必须分文档标注所有用到的原文信息。

\smallskip
\textbf{输入：}\\
相关文档内容：\texttt{\%s}

\smallskip
\textbf{输出格式（严格 JSON）：}

{\ttfamily
[\\
\quad \{\\
\qquad "question": "跨文档问题",\\
\qquad "answer": "跨文档推理答案",\\
\qquad "according": "文档1依据：xxx；文档2依据：xxx"\\
\quad \}\\
]
}

\\
\bottomrule
\end{tabular}

\caption{Prompt Template for Generating Cross-Document Multi-Hop Question--Answer Pairs.}
\label{tab:cross-document-multihop-prompt}

\end{table}